% !TEX root = task-4-12-3.tex
% Editable task bank. Append new BankTask blocks at the end of the document.
\documentclass[12pt,a4paper]{article}
\usepackage{physics-textbook}
\usepackage{graphicx,needspace}
\usetikzlibrary{decorations.pathmorphing}
\newif\ifShowSolutions
\ifdefined\TutorEdition
  \ShowSolutionstrue
\else
  \ShowSolutionsfalse % Change to \ShowSolutionstrue to display solutions.
\fi
\newcounter{banktask}
\newsavebox{\BankTaskBox}
\newcommand{\BankTaskContent}[3]{%
  \noindent\textbf{\thebanktask.}\quad #2\par
  \ifShowSolutions
    \Needspace{50mm}% Keep the solution heading with its opening text or diagram.
    #3\par
    \ifdefined\TutorEdition
      \ifcsname TutorSource#1\endcsname
        \nopagebreak\smallskip{\small\csname TutorSource#1\endcsname\par}%
      \fi
    \fi
  \fi
}
\newcommand{\BankTask}[3]{%
  \refstepcounter{banktask}\par\medskip
  \sbox{\BankTaskBox}{\begin{minipage}{\linewidth}%
    \BankTaskContent{#1}{#2}{#3}\end{minipage}}%
  \ifdim\dimexpr\ht\BankTaskBox+\dp\BankTaskBox\relax>\dimexpr\textheight-20mm\relax
    \Needspace{35mm}\BankTaskContent{#1}{#2}{#3}%
  \else
    \Needspace{\dimexpr\ht\BankTaskBox+\dp\BankTaskBox\relax}%
    \noindent\usebox{\BankTaskBox}\par
  \fi
  \medskip
}
\begin{document}
\BookSubsection{4.12.3}{Дифракционная решётка}


% Source 43126: https://phys-ege.sdamgia.ru/problem?id=43126
\BankTask{43126}{
К решетке вплотную поставили линзу, которая собирает лучи в одной точке на экране. Линза расположена между решеткой и экраном. Расстояние от линзы до экрана $2\,\text{м}$. Ширина спектра второго порядка $8\,\text{см}$. Длина волны красного света $8$ · $10^{-7}$ и фиолетового света $4$ · $10^{-7}$ м. Найдите, сколько штрихов приходится на $1\,\text{см}$ решетки. $\sin \varphi \approx \tan \varphi \approx \varphi.$
}{
\smallskip\noindent\textbf{Решение}\par
Условие дифракционного максимума $d \sin \varphi =k\lambda,$ где для малых углов $\sin \varphi \approx \tan \varphi = \frac{b}{L} .$
Тогда для максимума фиолетового и красного света получаем

$\frac{db_1}{L} =k\lambda_1$
$\frac{db_2}{L} =k\lambda_2.$

Откуда находим период решетки:

\[d= \frac{Lk ( \lambda_1 - \lambda_2 )}{\Delta b} = \frac{2 \cdot 2 \cdot ( 8 \cdot 10 ^{ - 7 } - 4 \cdot 10 ^{ - 7 } )}{8 \cdot 10 ^{ - 2 }} =2 \cdot 10 ^{ - 5 } \,\text{м}=2 \cdot 10 ^{ - 3 } \,\text{см}.\]

Тогда на $1\,\text{см}$ приходится количество штрихов:

\[N= \frac{1\,\text{см}}{d} = \frac{1}{2 \cdot 10 ^{ - 3 }} =500.\]

Ответ: $500$.
}

% Source 9333: https://phys-ege.sdamgia.ru/problem?id=9333
\BankTask{9333}{
Плоская монохроматическая световая волна с частотой $8{,}0$ · $10^{14}$ Гц падает по нормали на дифракционную решетку. Параллельно решетке позади нее размещена собирающая линза с фокусным расстоянием $21\,\text{см}$. Дифракционная картина наблюдается на экране в задней фокальной плоскости линзы. Расстояние между ее главными максимумами $1$-го и $2$-го порядков равно $18\,\text{мм}$. Найдите период решетки. Ответ выразите в микрометрах (мкм), округлив до десятых. Считать для малых углов ($\varphi \ll 1$ в радианах) $\tan \varphi \approx \sin \varphi \approx \varphi.$
}{
\smallskip\noindent\textbf{Решение}\par
Запишем условия максимума для $1$-го и $2$-го порядков:

\[\begin{cases}d\sin\varphi_1=\lambda,\\d\sin\varphi_2=2\lambda.\end{cases}\quad\Longleftrightarrow\quad\begin{cases}\tan\varphi_1\approx\sin\varphi_1=\frac{\lambda}{d}=\frac{c}{\nu d},\\\tan\varphi_2\approx\sin\varphi_2=\frac{2\lambda}{d}=\frac{2c}{\nu d}.\end{cases}\]

На экране максимумы смещены на расстояния $x_1=f \tan \varphi_1$ и $x_2=f \tan \varphi_2$ от центрального максимума, где f  — фокусное расстояние линзы. По условию $x_2 - x_1=\Delta x = 18\,\text{мм}.$ Из чего получаем

\[\frac{2fc}{\nu d} - \frac{fc}{\nu d} =\Delta x \Longleftrightarrow d= \frac{fc}{\nu \Delta x} = \frac{0{,}21 \cdot 3 \cdot 10 ^{8}}{8{,}0 \cdot 10 ^{ 14 } \cdot 18 \cdot 10 ^{ - 3 }} \approx 4{,}4 \cdot 10 ^{ - 6 } \,\text{м} = 4{,}4\,\text{мкм}.\]

Ответ: $4{,}4$
}

\end{document}
